\begin{workedexample}[Worked Example: Beat frequency and resolution]
A linear chirp of slope $S = B/T$ produces a beat frequency proportional to
range, $f_b = 2RS/c = 2RB/(cT)$. Range resolution depends only on sweep bandwidth:
$\Delta R = c/(2B)$. For $B = 150\ \text{MHz}$,
\[ \Delta R = \frac{3\times10^8}{2\times150\times10^6} = 1\ \text{m}. \]
Wider sweeps give finer range resolution; the beat frequency is read off an FFT of
the mixed signal.
\end{workedexample}

\begin{keytakeaways}
\begin{itemize}[leftmargin=1.4em,itemsep=2pt]
\item FMCW beat frequency $f_b = 2RS/c$ (slope $S = B/T$) encodes range.
\item Range resolution $\Delta R = c/(2B)$ depends only on sweep bandwidth.
\item A 2-D FFT over chirps yields the range-Doppler map (range and velocity).
\end{itemize}
\end{keytakeaways}

\begin{exercises}
\item Range resolution for $B = 1\ \text{GHz}$?
\item To halve $\Delta R$, how should $B$ change?
\item Does the beat frequency increase or decrease with range?
\end{exercises}

\furtherreading{Skolnik~\cite{skolnik}; Pozar~\cite{pozar}.}
